% Options for packages loaded elsewhere
\PassOptionsToPackage{unicode}{hyperref}
\PassOptionsToPackage{hyphens}{url}
\documentclass[
  11pt,
]{article}
\usepackage{xcolor}
\usepackage[margin=25mm]{geometry}
\usepackage{amsmath,amssymb}
\setcounter{secnumdepth}{-\maxdimen} % remove section numbering
\usepackage{iftex}
\ifPDFTeX
  \usepackage[T1]{fontenc}
  \usepackage[utf8]{inputenc}
  \usepackage{textcomp} % provide euro and other symbols
\else % if luatex or xetex
  \usepackage{unicode-math} % this also loads fontspec
  \defaultfontfeatures{Scale=MatchLowercase}
  \defaultfontfeatures[\rmfamily]{Ligatures=TeX,Scale=1}
\fi
\usepackage{lmodern}
\ifPDFTeX\else
  % xetex/luatex font selection
\fi
% Use upquote if available, for straight quotes in verbatim environments
\IfFileExists{upquote.sty}{\usepackage{upquote}}{}
\IfFileExists{microtype.sty}{% use microtype if available
  \usepackage[]{microtype}
  \UseMicrotypeSet[protrusion]{basicmath} % disable protrusion for tt fonts
}{}
\makeatletter
\@ifundefined{KOMAClassName}{% if non-KOMA class
  \IfFileExists{parskip.sty}{%
    \usepackage{parskip}
  }{% else
    \setlength{\parindent}{0pt}
    \setlength{\parskip}{6pt plus 2pt minus 1pt}}
}{% if KOMA class
  \KOMAoptions{parskip=half}}
\makeatother
\setlength{\emergencystretch}{3em} % prevent overfull lines
\providecommand{\tightlist}{%
  \setlength{\itemsep}{0pt}\setlength{\parskip}{0pt}}
\usepackage{bookmark}
\IfFileExists{xurl.sty}{\usepackage{xurl}}{} % add URL line breaks if available
\urlstyle{same}
% fallback for those not using the hyperref driver hyperxmp:
\makeatletter
\@ifundefined{xmpquote}{\newcommand{\xmpquote}[1]{#1}}{}
\makeatother
\hypersetup{
  pdftitle={Finite mean angular conditioning from a bounded semialgebraic graph},
  pdfauthor={Matthew J. Colbrook},
  hidelinks,
  pdfcreator={LaTeX via pandoc}}

\title{Finite mean angular conditioning from a bounded semialgebraic
graph}
\author{Matthew J. Colbrook}
\date{11 September 2026}

\usepackage{xurl}
\setlength{\emergencystretch}{3em}
\begin{document}
\maketitle

Department of Applied Mathematics and Theoretical Physics, University of
Cambridge, Cambridge, United Kingdom. Email: m.colbrook@damtp.cam.ac.uk.

This argument was developed with AI assistance at the author's request.
Independent review is recorded separately; no priority claim or external
human peer-review claim is made.

\subsection{1. Statement and model}\label{statement-and-model}

Fix a real tensor format with order \(d\ge3\), mode sizes \(n_j\ge2\),
and rank \(r\ge3\), subject to generic complex identifiability. Write
\(E\) for the ambient Euclidean tensor space with Frobenius norm. Let
\(M\) be the smooth identifiable real rank-\(r\) locus equipped with its
induced volume. As in the canonical TR-06 statement, the input density
is \[
d\mu(A)=Z^{-1}e^{-\|A\|^2/2}\,dV_M(A).
\] For any local ordering of the unique rank-one summands, let
\(\Psi(A)=(a_1,\ldots,a_r)\) and \[
f(A)=\left(\frac{a_1}{\|a_1\|},\ldots,\frac{a_r}{\|a_r\|}\right),
\qquad \kappa_{\rm ang}(A)=\|Df(A)\|_2.
\] The derivative uses the induced tangent norm on \(M\) and the product
Frobenius norm on \(E^r\). Changing the ordering is an isometry and does
not change this norm.

\textbf{Theorem.} For every format and rank under these hypotheses,
\(\mathbb E_\mu\kappa_{\rm ang}<\infty\).

This is the volume-Gaussian input model of Beltran, Breiding and
Vannieuwenhoven, Definition 1.3 in their preprint; their Conjecture 1.10
asks for this conclusion. It is not the distribution obtained by drawing
summands independently. We prove integrability, not a format-uniform
numerical upper bound or an efficient decomposition algorithm.

\subsection{2. Bounded graphs control a first
derivative}\label{bounded-graphs-control-a-first-derivative}

We use the standard finite-volume property of semialgebraic sets: a
bounded semialgebraic set of dimension at most \(s\) has finite
\(s\)-dimensional Hausdorff measure. A precise reference is Hardt,
Lambrechts, Turchin and Volic, Theorem 2.4. Closedness is not required.

Here is its application to a finite collection of local branches. Let
\(L\subset\mathbb R^a\) be a bounded smooth semialgebraic manifold of
dimension \(s\ge1\). Suppose \(G\subset L\times\mathbb R^b\) is bounded
and semialgebraic, its projection to \(L\) has finite nonempty fibers,
and locally it consists of smooth graphs of maps \(g\). Assume every
branch has the same derivative norm \(K(x)\), as happens when branches
differ by fixed coordinate permutations. Finite fibers imply
\(\dim G=s\). Thus \(\mathcal H^s(G)<\infty\).

For a branch the graph embedding \(x\mapsto(x,g(x))\) has volume
Jacobian \[
J_g(x)=\sqrt{\det(I+(Dg(x))^*Dg(x))}.
\] If \(\sigma_1,\ldots,\sigma_s\) are the singular values, including
zeros, then \[
J_g=\prod_{j=1}^s\sqrt{1+\sigma_j^2}\ge\max_j\sigma_j=K(x).
\] The ordinary area formula on each smooth graph therefore gives \[
\int_L K(x)\,dV_L(x)\le \mathcal H^s(G)<\infty. \tag{1}
\] For precision, take a countable local graph cover and a disjoint
measurable partition subordinate to its projection. On each piece count
every distinct branch once. The area formula sums the Jacobians of those
branches, and the sum is at least \(K\). The corresponding graph pieces
are disjoint. Monotone convergence proves (1), without assuming
integrability or a global ordering in advance. Lower-dimensional
semialgebraic pieces may be removed without affecting either integral.
Equivalently one may first use a finite smooth semialgebraic
stratification. Only local graph calculus and the cited bounded-volume
theorem are used here.

\subsection{3. The normalized decomposition
graph}\label{the-normalized-decomposition-graph}

The cone of nonzero rank-one tensors is a smooth semialgebraic manifold:
its algebraic equations are the flattening minors, with the origin
removed. The sum of \(r\) such tensors is a polynomial map on their
product. Generic complex identifiability implies that the ordered map
has finite generic fibers, consisting of permutations. In characteristic
zero it is unramified on a dense open set of its image. Restrict to that
regular identifiable locus and to real points. The omitted algebraic or
semialgebraic exceptional subset has smaller dimension and induced
volume zero. This is the regular, full-measure cone used in Proposition
2.2 of the source paper. We use local inverses and do not assume a
global labeling of summands.

Denote the resulting smooth semialgebraic cone by \(M^\circ\). Its
dimension is \[
k=r\left(1+\sum_{j=1}^d(n_j-1)\right)>1.
\] The dimension formula follows because the nonzero rank-one cone has
the displayed summand dimension and the ordered addition map is
generically finite. Let \[
L=M^\circ\cap\{A:\|A\|=1\}.
\] The radial direction is tangent to the cone, so the unit sphere is
transverse and \(L\) is smooth of dimension \(s=k-1\ge1\).

Define \(G\) to consist of all \[
\left(A,\frac{a_1}{\|a_1\|},\ldots,\frac{a_r}{\|a_r\|}\right),
\quad A\in L,\quad A=\sum_i a_i,
\] where \((a_1,\ldots,a_r)\) ranges over its ordered minimal
decompositions. The set is semialgebraic: the addition, rank-one and
normalization relations are polynomial equalities and inequalities after
introducing positive variables \(t_i\) with \(t_i^2=\|a_i\|^2\) and
\(t_i u_i=a_i\), and then projecting away \(a_i,t_i\). Identifiability
can likewise be expressed by a quantified polynomial formula, or imposed
by restriction to the regular identifiable semialgebraic locus above.

This graph is bounded, even though the unnormalized summands can
diverge: \(\|A\|=1\) and each normalized summand has norm one. Its
projection fibers are finite and nonempty. No two summands in a minimal
decomposition can be proportional, since they could be combined or
canceled to reduce the number of terms. Consequently the \(r!\)
permutations yield distinct ordered normalized tuples. Near a regular
point the graph consists of their smooth local branches. The derivative
norms are identical under these permutations.

Apply (1) to the restriction of \(f\) to the link. It follows that \[
I_L:=\int_L\|D(f|_L)(A)\|_2\,dV_L(A)<\infty. \tag{2}
\] This argument is unaffected by a branch permutation around a loop. It
makes no claim that the ordered addition map is globally one-to-one.

\subsection{4. Radial integration}\label{radial-integration}

For \(t>0\), uniqueness gives \(\Psi(tA)=t\Psi(A)\) locally with
compatible labels. Therefore \(f(tA)=f(A)\). The radial derivative
vanishes. The orthogonal decomposition \[
T_A M^\circ=\mathbb R A\oplus T_A L\qquad(A\in L)
\] shows \(\|Df(A)\|_2=\|D(f|_L)(A)\|_2\) and, at general radius, \[
\kappa_{\rm ang}(tA)=t^{-1}\|D(f|_L)(A)\|_2. \tag{3}
\] The polar map \((t,A)\mapsto tA\) on the cone has volume element
\(t^{k-1}\,dt\,dV_L(A)\). Tonelli's theorem, (2) and (3) give \[
\int_{M^\circ}e^{-\|A\|^2/2}\kappa_{\rm ang}(A)\,dV_M(A)
= I_L\int_0^\infty t^{k-2}e^{-t^2/2}\,dt<\infty.
\] The radial integral converges at zero because \(k>1\) and at infinity
because of Gaussian decay. Similarly \[
Z=\mathcal H^{k-1}(L)\int_0^\infty t^{k-1}e^{-t^2/2}\,dt
\] is finite and positive: \(L\) is nonempty and smooth of its stated
dimension and is bounded semialgebraic. Restoring the removed null set
does not change these integrals. Dividing by \(Z\) proves the theorem.

The proof also covers rank two under the same assumptions, consistently
with the known theorem. It does not imply integrability of the ordinary
condition number: the graph of the unnormalized summands need not be
bounded. Nor does it prove higher moments, since a graph Jacobian need
not dominate an arbitrary power of its largest singular value.

\subsection{References}\label{references}

\begin{enumerate}
\def\labelenumi{\arabic{enumi}.}
\tightlist
\item
  C. Beltran, P. Breiding and N. Vannieuwenhoven, \emph{The Average
  Condition Number of Most Tensor Rank Decomposition Problems Is
  Infinite}, Foundations of Computational Mathematics 23 (2023),
  433-491. \href{https://arxiv.org/abs/1903.05527}{Preprint}, Definition
  1.3, Conjecture 1.10, Proposition 2.2 and Lemma 6.1;
  \href{https://doi.org/10.1007/s10208-022-09551-1}{journal}, equation
  (6), Theorem 3 and Conjecture 2.
\item
  R. Hardt, P. Lambrechts, V. Turchin and I. Volic, \emph{Real homotopy
  theory of semi-algebraic sets}, Algebraic \& Geometric Topology 11
  (2011), 2477-2545, Theorem 2.4.
  \href{https://msp.org/agt/2011/11-5/agt-v11-n5-p01-s.pdf}{Primary
  text}.
\item
  \href{https://github.com/ajt60gaibb/OpenProblemsInNLA/blob/main/tensor-computations/TR-06/README.md}{Canonical
  TR-06 statement}, checked 11 September 2026. The bounded search and
  claim inventory are recorded separately in the submission record.
\end{enumerate}

\end{document}
